% 01-tarifname.tex

\begin{center}
    \textbf{TARİFNAME} \\
    \vspace{1em}
    \textbf{BULUŞUN BAŞLIĞINI BURAYA YAZINIZ}
\end{center}

\section*{Teknik Alan}
% Buluşun hangi teknik alanla ilgili olduğu bu kısımda kısaca belirtilir.
Bu buluş, ... ile ilgilidir.

\section*{Önceki Teknik}
% Buluşun anlaşılması için yurt içi ve yurt dışındaki benzer buluş, teknoloji ya da ürünler burada anlatılır.

\section*{Buluşun Amacı}
% Buluşun tekniğin bilinen durumuna kıyasla getirdiği yenilikler ve çözdüğü problemler anlatılmalıdır.

\section*{Şekillerin Açıklaması}
% Başvuruda resim bulunuyorsa, bu resimlerdeki şekillerin kısa açıklamaları liste halinde verilir.
% Örnek: Şekil 1: Buluş konusu ...'nın perspektif görünüşüdür.

\section*{Şekillerdeki Referansların Açıklaması}
% Şekillerde bulunan parçalar vb. referans işaretleri ile gösterilir.
% Örnek:
% 1: Gövde
% 2: Kapak

\section*{Buluşun Açıklaması}
% Buluşun tüm detayı ile anlatıldığı kısımdır.

\section*{Buluşun Sanayiye Uygulanma Biçimi}
% Buluşun nerede kullanılabileceği ve sanayiye nasıl uygulanacağı belirtilir.
